% ===================================================================
%  도표 제 2 호 — 사람의 몸. 쉬는 시간. 놀이.
%
%  Traduction, faite en 2026, en coreen d'aujourd'hui, sur l'ido de
%  texto/io. Regles de la colonne : voir 10-dopyo-01.tex.
%
%  L'ANATOMIE EST LE MEILLEUR TEST QU'UNE LANGUE PUISSE SUBIR, parce
%  qu'elle nomme cinquante choses d'affilee sans une seule
%  periphrase possible. Le coreen les a toutes, et pour la plupart
%  d'un mot de fond : 이마 le front, 관자놀이 la tempe, 속눈썹 les
%  cils, 콧구멍 les narines, 구레나룻 les favoris, 목덜미 la nuque,
%  갈비뼈 les cotes, 옆구리 le flanc, 팔꿈치 le coude, 손목 le
%  poignet, 손톱 les ongles, 넓적다리 la cuisse, 오금 le jarret,
%  장딴지 le mollet, 복사뼈 la malleole, 발바닥 la plante, 발꿈치 le
%  talon. Pas un emprunt dans les cinquante.
%
%  LES JEUX, EUX, SE PARTAGENT EN TROIS, ET LA COUPURE EST NETTE.
%
%   * Ceux que la Coree joue et nomme : 굴렁쇠 le cerceau (14), 팽이
%     la toupie (18), 구슬치기 les billes (21), 죽마 les echasses
%     (24), 숨바꼭질 cache-cache (25), 까막잡기 le colin-maillard
%     (31), 연 le cerf-volant (33), 그네 la balancoire (35). Huit
%     jeux, huit mots coreens, et aucun n'est un calque.
%   * Ceux qui sont venus avec leur nom : 테니스 (34), 크로케 (15),
%     빌보케 (19) — le bilboquet n'a pas de nom coreen, et lui en
%     fabriquer un ferait croire a un jeu qu'on connait.
%   * Ceux qu'il faut decrire, faute de nom : le jeu de barres (38),
%     la main chaude (28), le jeu de l'ours (32). La tentation etait
%     forte pour le premier : 진놀이 est un jeu coreen de deux camps
%     qui se poursuivent, exactement la meme forme. On ne l'ecrit
%     pas. Ce serait remplacer le jeu du livret par un autre, et la
%     regle est la meme depuis le tableau 5 telougou — on modernise
%     la langue, jamais les choses.
%
%  UN SEUL RETOURNEMENT SUR LES SEPT PAIRES DE parigi.py :
%  « kegli (16) quin li abatas per bulo (17) ». Les six autres sont
%  des frontieres de phrase et des coordinations. Le tableau est
%  presque tout entier une enumeration — le corps, puis les jeux —
%  et l'enumeration ne s'inverse pas.
% ===================================================================

%%K t02-apar-1 apar
\VUsaut{4.0mm}
\VUtitre{15.0pt}{60}{도표 제 2 호}
\VUsaut{2.0mm}
\VUtitre{11.4pt}{0}{사람의 몸. --- 쉬는 시간.}
\VUtitre{11.4pt}{0}{놀이.}

%%K t02-tit-0 sub
\VUtitre{13.2pt}{0}{사람의 몸.}
\VUsaut{2.4mm}
\VUcentre{10.2pt}{0}{\textit{(자연 과학의 간단한 한 과.)}}
\VUsaut{3.0mm}

%%K t02-tit-1 sub pos=t02-01-1
\VUpk{10.2pt}{40}{머리.}
\VUsaut{3.0mm}

%%K t02-01-1 p
1. --- 사람 몸의 주된 부분은 \VUgras{머리}, \VUgras{몸통}, \VUgras{팔다리}다.

%%K t02-02-1 p
2. --- 머리는 두 부분으로 나뉜다. \VUgras{두개골}은 \VUgras{뇌}를 담고 있고
\VUgras{머리털}이 그것을 덮으며, 또 하나는 \VUgras{얼굴}이다.

%%K t02-03-1 p
3. --- 우리 그림에서는 두개골도 뇌도 보이지 않는다. 그러나 얼굴의 중요한 부분들은 아주
또렷하게 보인다.

%%K t02-04-1 p
4. --- 먼저 \VUgras{이마}가 있다. 그것은 힘 있는 지성과 늘 움직이는 생각을 드러낸다.
\VUgras{관자놀이}는 아주 트여 있다. \VUgras{눈}은 생기 있고 크게 열려 있다. 짙은
\VUgras{눈썹}과 긴 \VUgras{속눈썹}이 그것을 가려 준다.

%%K t02-05-1 p
5. --- 다음으로 \VUgras{코}와 \VUgras{콧구멍}이 있다. 코는 냄새를 맡고 숨을 쉬는 데
쓰인다. 코 양쪽에는 \VUgras{뺨}이, 더 멀리에는 \VUgras{귀}가 있다. 얼굴 아래쪽은
\VUgras{입}과 \VUgras{턱}으로 이루어진다.

%%K t02-06-1 p
6. --- 그것들은 잘 보이지 않는다. \VUgras{콧수염}과 \VUgras{구레나룻}이 우리에게서
그것을 가리기 때문이다. 그러나 우리는 입에 두 \VUgras{입술}과 \VUgras{위턱},
\VUgras{아래턱}, 그리고 그 \VUgras{이}와 \VUgras{혀}가 있다는 것을 안다. 입으로 우리는
말하고 먹는다. 혀로 우리는 음식을 맛본다.

%%K t02-07-1 p
7. --- 눈, 귀, 코, 입은 손가락과 마찬가지로 다섯 감각의 기관이다. 곧 \VUgras{시각},
\VUgras{청각}, \VUgras{후각}, \VUgras{미각}, \VUgras{촉각}이다.

%%K t02-08-1 p
8. --- 그러니 머리는 사람 몸에서 가장 흥미로운 부분 가운데 하나다.

%%K t02-tit-2 sub
\VUsaut{4.4mm}
\VUpk{10.2pt}{40}{몸통.}
\VUsaut{3.0mm}

%%K t02-09-1 p
9. --- \VUgras{목}이 머리를 몸통에 잇는다. 거기에는 \VUgras{목구멍}과 \VUgras{목덜미}가
들어 있다.

%%K t02-10-1 p
10. --- 몸통에는 없어서는 안 될 기관도 여럿 들어 있다. \VUgras{가슴}에는 우리가 숨쉬게
해 주는 \VUgras{허파}와, 삶의 중심인 \VUgras{심장}이 있다. 심장이 뛰기를 멈추면 살아
있는 것은 죽는다.

%%K t02-11-1 p
11. --- \VUgras{갈비뼈}가 가슴을 지탱한다.

%%K t02-12-1 p
12. --- 몸통의 나머지 부분은 \VUgras{옆구리}, \VUgras{등}, \VUgras{어깨},
\VUgras{배}다. 우리는 잠자려고 옆구리나 등을 대고 눕고, 짐은 어깨에 진다.

%%K t02-13-1 p
13. --- 배 속에는 \VUgras{위}와 \VUgras{창자}가 있다. 그것들은 음식을 삭이는 데 쓰인다.

%%K t02-tit-3 sub
\VUsaut{4.4mm}
\VUpk{10.2pt}{40}{팔다리.}
\VUsaut{3.0mm}

%%K t02-14-1 p
14. --- 우리에게는 \VUgras{팔다리}가 넷 있다. \VUgras{팔} 둘과 \VUgras{다리} 둘이다.
팔로 우리는 물건을 잡고 들고 올리고 거두어들이며, 다리로 우리는 서고 걷고 달린다.

%%K t02-15-1 p
15. --- \VUgras{어깨}가 팔을 몸에 잇는다. 아주 튼튼한 \VUgras{근육}인 이두박근이 팔을
움직이고, \VUgras{팔꿈치}가 그 팔을 \VUgras{아래팔}에 잇는다. \VUgras{손목}은
\VUgras{손}의 관절을 이루고, 손은 \VUgras{손가락}과 \VUgras{손톱}으로 끝난다. 손등에서
우리는 \VUgras{정맥}(과 동맥)을 알아볼 수 있는데, 그 안으로 \VUgras{피}가 흐른다.

%%K t02-16-1 p
16. --- 다리의 부분은 이러하다. \VUgras{넓적다리}, \VUgras{무릎}과 \VUgras{오금},
\VUgras{장딴지} — 그 \VUgras{살}은 \VUgras{살갗}으로 덮여 있다 —, \VUgras{복사뼈},
\VUgras{발}. 다리의 \VUgras{뼈}는 아주 굵고 아주 튼튼하다. 발 끝에는
\VUgras{발가락}이 있다. 나머지 부분은 \VUgras{발바닥}과 \VUgras{발꿈치}다.

%%K t02-17-1 p
17. --- 사람 몸에 대한 거의 온전한 설명이 이러하다.

%%K t02-17-2 p
\textit{알림.} ---
\textit{« 나는 ...(머리 따위)가 아프다 »라는 말투의 도움으로, 선생님은 이 낱말들 전부를
좋거나 나쁜} \VUgras{몸 상태} \textit{의 관점에서 다시 쓰게 할 수 있을 것이다.}

%%K t02-tit-4 sub
\VUsaut{5.0mm}
\VUpk{10.2pt}{40}{체조.}
\VUsaut{2.4mm}
\VUcentre{10.2pt}{0}{\textit{(학생 하나의 이야기.)}}
\VUsaut{3.0mm}

%%K t02-18-1 p
18. --- 몸이 부드럽고 날래고 튼튼하려면 우리는 다리와 팔을 단련해야 한다. 그래서 우리는
놀고 \VUgras{체조}를 한다.

%%K t02-19-1 p
19. --- 주중에는 이 두 가지를 고등학교 운동장에서 한다(도표 제 1 호를 보라). 그러나
목요일과 일요일에는 넓은 잔디밭으로 가는데, 거기에는 달리고 놀 자리가 있다.

%%K t02-20-1 p
20. --- 선생님들과 부모님들이 이따금 우리와 함께 그곳에 가신다. 그러나 운동을 이끄시는
분은 \VUgras{체조 선생님}\textsuperscript{(12)}이다.

%%K t02-21-1 p
21. --- 잔디밭 들머리 왼쪽에는 \VUgras{체조 기구}\textsuperscript{(1)}가 있다. 우리는
\VUgras{사다리}\textsuperscript{(2)}를 타고 오르고, \VUgras{링}\textsuperscript{(3)}과
\VUgras{공중그네}\textsuperscript{(4)}에서 거꾸로 돌며, \VUgras{줄사다리}\textsuperscript{(5)}나
\VUgras{매듭 밧줄}\textsuperscript{(6)} 꼭대기까지 손힘으로 올라간다.

%%K t02-22-1 p
22. --- 그동안 우리 동무들은 \VUgras{목마}\textsuperscript{(7)}나
\VUgras{평행봉}\textsuperscript{(11)}을 뛰어넘는다. 어떤 아이들은
\VUgras{권투를 하고}\textsuperscript{(8)}, 어떤 아이들은 얼굴 가리개를 쓰고
\VUgras{플뢰레}로 \VUgras{펜싱을 한다}\textsuperscript{(9)}. 또 어떤 아이들은
\VUgras{역기를 든다}\textsuperscript{(10)}.

%%K t02-tit-5 sub
\VUsaut{4.6mm}
\VUpk{10.2pt}{40}{놀이.}
\VUsaut{3.0mm}

%%K t02-23-1 p
23. --- 체조하는 아이들 둘레에서 사내아이와 계집아이들이 여러 가지 \VUgras{놀이}를 한다.
계집아이들은 자기 \VUgras{굴렁쇠}\textsuperscript{(14)}로 논다. 그 아이들은
\VUgras{줄}\textsuperscript{(13)}을 넘거나, 오라비와 사촌들을
\VUgras{크로케}\textsuperscript{(15)} 한 판에 부른다. 상대가 작은 문을 쉽게 지나가겠다고
생각하는 바로 그때, 자기 \VUgras{나무망치}로 그 상대의 \VUgras{공}을 밀어내는 것이 그
아이들에게는 신난다!

%%K t02-24-1 p
24. --- 사내아이들은 더 힘쓰는 놀이를 좋아한다. 그들은 즐겨
\VUgras{나무 핀}\textsuperscript{(16)} 놀이를 하는데, 그것을
\VUgras{공}\textsuperscript{(17)}으로 솜씨 좋게 쓰러뜨린다.
\VUgras{팽이}\textsuperscript{(18)}나 \VUgras{빌보케}\textsuperscript{(19)}, 심지어
\VUgras{개구리 놀이}\textsuperscript{(20)}도 마다하지 않는다. 그러나 그들이 가장 좋아하는
놀이는 역시 \VUgras{구슬치기}\textsuperscript{(21)}와
\VUgras{벽치기 공놀이}\textsuperscript{(22)}다. \VUgras{온실}\textsuperscript{(26)}의 벽이
가벼운 공을 되튀기는 데 아주 알맞기 때문이다.

%%K t02-25-1 p
25. --- 어린 요안네스는 자기 \VUgras{죽마}\textsuperscript{(24)}를 타고 다니는데, 아주
솜씨가 좋고 균형을 잘 지킬 줄 안다. 그 곁에서는 동무 몇이 그 온실 안과 \VUgras{산울타리}
뒤에서 \VUgras{숨바꼭질}\textsuperscript{(25)}을 한다. 다른 아이들은
\VUgras{때린 사람 알아맞히기}\textsuperscript{(28)}나
\VUgras{날개 공}\textsuperscript{(29)}을 더 좋아하는데, 그 공은 한 \VUgras{라켓}에서 다른
라켓으로 날아다닌다.

%%K t02-noto-1 noto
\VUnotes{143.2mm}{%
(*) 아이들 놀이에는 온전히 나라마다 다른 이름이 붙는 일이 흔하다. 우리는 그것을 이도로
되도록 잘 부르려고 애썼다. 때로는 그 놀이를 특징짓는 동작 자체에서 실마리를 얻고, 때로는
이 나라 저 나라의 이름이 지나치게 치우치지 않을 때 그 이름을 골랐다. 보기를 들면
\VUgras{통 놀이}(독일과 프랑스)는 \VUgras{개구리 놀이} \textsuperscript{(20)}(영국)인데,
거기에서 노는 이들은 구멍 뚫린 통 위에 입을 벌리고 선 쇠개구리의 입 안으로 자기 둥근
쇠붙이를 던져 넣으려고 애쓴다.
\VUgras{목말 뛰어넘기}\textsuperscript{(30)}가 \VUgras{등 짚고 뛰어넘기}와 다른 점은 이것뿐이다.
머리 위로 뛰어넘으려고 노는 이들이 짝의 어깨에 손을 짚는다는 것. 「\VUgras{등 짚고
뛰어넘기}」에서는 그 등에만 손을 짚는다. --- 또는 참가자 가운데 몇이 몸을 굽히고 있는 동안 다른 이들이 어깨에 손을 짚고 그 등
위로 뛰어넘는데, 앞의 이들은 그 충격을 굽히지 않고 견뎌야 하고 뒤의 이들은 균형을 잃지
않고 뛰어야 한다(도표 자체를 보라).%
}

%%K t02-26-1 p
26. --- 잔디밭 한가운데에는 나무가 둘 있다. 그 가운데 하나의 밑동에서 사내아이 일고여덟이
\VUgras{목말 뛰어넘기}\textsuperscript{(30)} 한 판을 벌였다 (*). 이 놀이는 어느 나라에서나
알려져 있지는 않다. 그러나 \VUgras{등 짚고 뛰어넘기}가 무엇인지는 모두가 안다. 다만
아이들은 이번에 그것을 하지는 않는다.

%%K t02-tit-6 sub
\VUsaut{5.0mm}
\VUpk{10.2pt}{40}{바깥에서 하는 놀이.}
\VUsaut{3.4mm}

%%K t02-27-1 p
27. \VUgras{까막잡기}\textsuperscript{(31)}도 아주 재미있다. 계집아이와 사내아이가 번갈아
\VUgras{눈가리개}를 눈에 대고, 잡히도록 내버려 두는 서투른 아이를 붙든다.

%%K t02-28-1 p
28. --- 잔디밭 한가운데, 아주 프랑스다우면서도 좀 거친 놀이가 있다. 그것은
\VUgras{곰놀이}\textsuperscript{(32)}인데, \VUgras{연}\textsuperscript{(33)}의 그토록
조용한 놀이보다, 또 영국 놀이인 \VUgras{테니스}\textsuperscript{(34)}보다도 훨씬 시끄럽다.

%%K t02-29-1 p
29. --- 그 마지막 운동은 큰 학생들의 기쁨이다. 우리 선생님들까지도 즐겨 하신다. 그것은
날램과 함께 큰 솜씨를 요구하고, 나에게는 아주 마음에 든다.
\VUgras{그네}\textsuperscript{(35)} 놀이는 계집아이들에게 맡긴다.

%%K t02-30-1 p
30. --- 자칫 위험해질 수도 있는 다른 영국 놀이는 나는 좋아하지 않는다.

%%K t02-30-1b p
내가 말하려는 것은 \VUgras{축구}\textsuperscript{(36)}다. 그것은 우리 형을 즐겁게 한다.
나는 우리의 오래된 \VUgras{바 놀이}\textsuperscript{(38)}가 더 좋다. 그것도 그만큼
몸에 이롭고 정말로 덜 거칠다.

%%K t02-tit-7 sub
\VUsaut{4.6mm}
\VUpk{10.2pt}{40}{자전거 타기와 공놀이.}
\VUsaut{3.0mm}

%%K t02-31-1 p
31. --- \VUgras{자전거}\textsuperscript{(37)}로 잔디밭을 도는 것도 나는 즐겨 한다.

%%K t02-32-1 p
32. --- 다음 해에는 \VUgras{승마}\textsuperscript{(39)}를 배울 것이다. 곱게 걷거나
빨리 달리고, 달음질치며 장애물을 뛰어넘는 말은 어찌나 아름다운지!

%%K t02-33-1 p
33. --- 여름에는 \VUgras{수영}\textsuperscript{(40)}을 배운다. 우리는 강의 맑고 시원한
물에서 즐겁게 자맥질하고 미역을 감는다. 이따금 마지막으로 종이
\VUgras{풍선}\textsuperscript{(41)}을 띄우는데, 더운 공기로 채우자마자 그것은 하늘로
의젓하게 올라간다.

%%K t02-34-1 p
34. --- 다른 놀이도 아주 많지만 다 헤아리지는 못하겠다. 이 간추림만 보아도 목요일마다
우리 고운 잔디밭에서 지루할 일이 없다는 것을 알 수 있다.

%%K t02-34-2 p
덧붙임. --- \textit{선생님은 이 장을 쉽게 채워 넣을 수 있을 것이다. 가장 중요한 놀이
하나하나를 더 자세히 풀어 주면서.}
